\section{Assumptions and notation}\label{sec:assume}

\subsection{Assumptions}

\begin{enumerate}[leftmargin=1.5em,itemsep=1pt]
\item \textbf{Station records represent their cities.} One GHCN station represents one city; Chinese sites are urban or airport stations, so they describe the built-up area rather than the whole prefecture. The Zhengzhou station's three-day total is 70\% of the official areal 449\,mm; station data drive every statistic and design depth, and the official figure enters once, as the scale anchor of the damage function, with the consequence of using the station total instead reported in Section~\ref{sec:eval}.
\item \textbf{Valid years are homogeneous years.} A year enters only with at least 300 valid precipitation days; partial years are discarded rather than rescaled, because rescaling would fabricate a seasonal distribution.
\item \textbf{The baseline is stationary.} The hazard model is stationary with an optional linear trend in location, supported by the failure of the Pettitt test at Zhengzhou \cite{pettitt1979} and by the scarcity of significant total-rainfall trends nationally.
\item \textbf{Precipitation is the only hazard driver.} Surge, riverine flooding and wind are not modelled; the documented damage in both events was pluvial.
\item \textbf{SCS-CN is a screening tool here.} Applied to a 1\,180\,km$^2$ city it estimates runoff volume, not street depth; CN is bracketed between 75 and 95 and the range is reported.
\item \textbf{Official losses are authoritative and comparable.} The three loss figures come from reports compiled under one national statistical standard; no inflation or reporting correction is applied.
\item \textbf{Damage grows as a power law in overflow,} with exponent 1.5 fixed by the observed convexity of urban flood damage and calibrated on one event.
\item \textbf{Costs are planning-grade unit costs,} varied in the sensitivity analysis rather than presented as tendered prices.
\end{enumerate}

The notation is collected in Table~\ref{tab:notation} in Appendix~\ref{app:notation}.